\subsubsection{General Output-Size Formula}\label{sec:general-output-size-formula}

Now we can extend the formula introduced conceptually in the padding section (\autopageref{eq:output-size-convolution-stride-1}). With padding but stride $1$, the output size is:
\begin{equation*}
    N_{\text{out}} = N_{\text{in}} + 2P - K + 1
\end{equation*}
Because the number of valid kernel starting positions was $N_{\text{in}} + 2P - K + 1$. Now we introduce stride $S$. We are no longer visiting \textbf{every} valid starting position, but only every $S$-th position. So the available travel distance is still:
\begin{equation*}
    N_{\text{in}} + 2P - K
\end{equation*}
If we move by steps of size $S$, the number of jumps that fit in that distance is:
\begin{equation*}
    N_{\text{out}} = \left\lfloor \frac{N_{\text{in}} + 2P - K}{S} \right\rfloor
\end{equation*}
Then we add $1$ because the initial position counts too:
\begin{equation}\label{eq:output-size-convolution-stride-general}
    N_{\text{out}} = \left\lfloor \frac{N_{\text{in}} + 2P - K}{S} \right\rfloor + 1
\end{equation}
\hl{The floor function is necessary because the last jump may not fit exactly} (e.g., if $N_{\text{in}}=8, K=3, P=0, S=2$, the jump to perform are $(8-3) / 2 = 2.5$, but we cannot place a kernel at half a stride, so we take the floor and add $1$ for the initial position, giving $3$ valid positions: $0, 2, 4$).

\highspace
\begin{flushleft}
    \textcolor{Green3}{\faIcon{link} \textbf{Relationship with the previous formulas}}
\end{flushleft}
If stride $S=1$, the general formula becomes:
\begin{equation*}
    N_{\text{out}} = \left\lfloor \frac{N_{\text{in}} + 2P - K}{1} \right\rfloor + 1 = N_{\text{in}} + 2P - K + 1
\end{equation*}
Which is exactly the formula we derived in the padding section (\autopageref{eq:output-size-convolution-stride-1}). If also there is no padding ($P=0$), we get simply:
\begin{equation*}
    N_{\text{out}} = \left\lfloor \frac{N_{\text{in}} - K}{S} \right\rfloor + 1 = \left\lfloor \frac{N_{\text{in}} - K}{1} \right\rfloor + 1 = N_{\text{in}} - K + 1
\end{equation*}

\highspace
\textcolor{Red2}{\faIcon{exclamation-triangle}} Finally, just like padding, stride affects \textbf{where the filter is evaluated}, not the filter weights themselves. \hl{Stride does not change the number of parameters.}